\documentclass[12pt]{article}
\usepackage{latexblog}

% ============================================================
% Blog Metadata
% ============================================================
\blogtitle{MySQL replication}
\blogdate{2019-01-21}
\blogtags{MySQL, 闲话编程, 存储}
\bloglang{zh}
\blogsummary{}

\begin{document}
\makeblogtitle

MySQL设置Replication后，可以支持Master库上的修改自动同步到Slave库上。利用Docker可以在本机尝试这种特性。

\section{配置Master}\label{ux914dux7f6emaster}

首先需要创建几个文件夹（略），用来挂载配置文件和数据。我们首先来配置Master库：

\begin{Shaded}
\begin{Highlighting}[]
\CommentTok{\# master/cnf/my.cnf}
\KeywordTok{[mysqld]}

\DataTypeTok{server{-}id}\OtherTok{=}\DecValTok{1}
\DataTypeTok{log{-}bin}\OtherTok{=}\StringTok{/var/lib/mysql/mysql{-}bin.log}
\DataTypeTok{binlog\_format}\OtherTok{=}\StringTok{MIXED}
\DataTypeTok{expire\_logs\_days}\OtherTok{=}\DecValTok{7}
\DataTypeTok{max\_binlog\_size}\OtherTok{=}\StringTok{50m}
\DataTypeTok{max\_binlog\_cache\_size}\OtherTok{=}\StringTok{256m}
\end{Highlighting}
\end{Shaded}

启动Master：

\begin{Shaded}
\begin{Highlighting}[]
\ExtensionTok{docker}\NormalTok{ run }\AttributeTok{{-}{-}name}\NormalTok{ mysql\_master }\DataTypeTok{\textbackslash{}}
    \AttributeTok{{-}{-}mount}\NormalTok{ type=bind,src=/Users/hfli/mysql{-}replication/master/cnf/my.cnf,dst=/etc/my.cnf }\DataTypeTok{\textbackslash{}}
    \AttributeTok{{-}{-}mount}\NormalTok{ type=bind,src=/Users/hfli/mysql{-}replication/master/data/,dst=/var/lib/mysql }\DataTypeTok{\textbackslash{}}
    \AttributeTok{{-}e}\NormalTok{ MYSQL\_ROOT\_PASSWORD=1125482715 }\DataTypeTok{\textbackslash{}}
    \AttributeTok{{-}d}\NormalTok{ mysql:5.7.24}
\end{Highlighting}
\end{Shaded}

然后需要登录到MySQL创建一个用来复制的用户

\begin{Shaded}
\begin{Highlighting}[]
\NormalTok{create user \textquotesingle{}replication\textquotesingle{} identified by \textquotesingle{}1153687060\textquotesingle{};}
\NormalTok{grant replication slave on *.* to \textquotesingle{}replication\textquotesingle{}@\textquotesingle{}\%\textquotesingle{} identified by \textquotesingle{}1153687060\textquotesingle{};}
\end{Highlighting}
\end{Shaded}

接下来需要看一下Master库的状态:

\begin{verbatim}
mysql> show master status;
+------------------+----------+--------------+------------------+-------------------+
| File             | Position | Binlog_Do_DB | Binlog_Ignore_DB | Executed_Gtid_Set |
+------------------+----------+--------------+------------------+-------------------+
| mysql-bin.000003 |      696 |              |                  |                   |
+------------------+----------+--------------+------------------+-------------------+
1 row in set (0.00 sec)
\end{verbatim}

\section{从库配置}\label{ux4eceux5e93ux914dux7f6e}

\begin{Shaded}
\begin{Highlighting}[]
\CommentTok{\# slave1/cnf/my.cnf}
\KeywordTok{[mysqld]}

\DataTypeTok{server{-}id}\OtherTok{=}\DecValTok{2}
\end{Highlighting}
\end{Shaded}

启动docker：

\begin{Shaded}
\begin{Highlighting}[]
\ExtensionTok{docker}\NormalTok{ run }\AttributeTok{{-}{-}name}\NormalTok{ mysql\_slave1 }\DataTypeTok{\textbackslash{}}
    \AttributeTok{{-}{-}mount}\NormalTok{ type=bind,src=/Users/hfli/mysql{-}replication/slave1/cnf/my.cnf,dst=/etc/my.cnf }\DataTypeTok{\textbackslash{}}
    \AttributeTok{{-}{-}mount}\NormalTok{ type=bind,src=/Users/hfli/mysql{-}replication/slave1/data/,dst=/var/lib/mysql }\DataTypeTok{\textbackslash{}}
    \AttributeTok{{-}{-}link}\NormalTok{ mysql\_master }\DataTypeTok{\textbackslash{}}
    \AttributeTok{{-}e}\NormalTok{ MYSQL\_ROOT\_PASSWORD=1125482715 }\DataTypeTok{\textbackslash{}}
    \AttributeTok{{-}d}\NormalTok{ mysql:5.7.24}
\end{Highlighting}
\end{Shaded}

然后即可启动Replication:

\begin{verbatim}
mysql> change master to master_host='mysql_master',master_user='replication',master_password='1153687060',master_log_file='mysql-bin.000003',master_log_pos=696;
Query OK, 0 rows affected, 2 warnings (0.03 sec)

mysql> start slave;
Query OK, 0 rows affected (0.00 sec)

mysql> show slave status\G;
*************************** 1. row ***************************
               Slave_IO_State: Waiting for master to send event
                  Master_Host: mysql_master
                  Master_User: replication
                  Master_Port: 3306
                Connect_Retry: 60
              Master_Log_File: mysql-bin.000003
          Read_Master_Log_Pos: 696
               Relay_Log_File: c17b953fb671-relay-bin.000002
                Relay_Log_Pos: 320
        Relay_Master_Log_File: mysql-bin.000003
             Slave_IO_Running: Yes
            Slave_SQL_Running: Yes
              Replicate_Do_DB:
          Replicate_Ignore_DB:
           Replicate_Do_Table:
       Replicate_Ignore_Table:
      Replicate_Wild_Do_Table:
  Replicate_Wild_Ignore_Table:
                   Last_Errno: 0
                   Last_Error:
                 Skip_Counter: 0
          Exec_Master_Log_Pos: 696
              Relay_Log_Space: 534
              Until_Condition: None
               Until_Log_File:
                Until_Log_Pos: 0
           Master_SSL_Allowed: No
           Master_SSL_CA_File:
           Master_SSL_CA_Path:
              Master_SSL_Cert:
            Master_SSL_Cipher:
               Master_SSL_Key:
        Seconds_Behind_Master: 0
Master_SSL_Verify_Server_Cert: No
                Last_IO_Errno: 0
                Last_IO_Error:
               Last_SQL_Errno: 0
               Last_SQL_Error:
  Replicate_Ignore_Server_Ids:
             Master_Server_Id: 1
                  Master_UUID: 83a0a667-1d50-11e9-b754-0242ac110002
             Master_Info_File: /var/lib/mysql/master.info
                    SQL_Delay: 0
          SQL_Remaining_Delay: NULL
      Slave_SQL_Running_State: Slave has read all relay log; waiting for more updates
           Master_Retry_Count: 86400
                  Master_Bind:
      Last_IO_Error_Timestamp:
     Last_SQL_Error_Timestamp:
               Master_SSL_Crl:
           Master_SSL_Crlpath:
           Retrieved_Gtid_Set:
            Executed_Gtid_Set:
                Auto_Position: 0
         Replicate_Rewrite_DB:
                 Channel_Name:
           Master_TLS_Version:
1 row in set (0.00 sec)
\end{verbatim}

一个有趣的问题：如果我修改了从库，会产生什么影响？例如已经有重复的数据，那么同步的时候就会报错，我们通过Last\_Error可以看到错误。

\begin{verbatim}
 Last_Errno: 1062
 Last_Error: Error 'Duplicate entry '2' for key 'PRIMARY'' on query. Default database: 'foo'. Query: 'INSERT INTO `foo`.`bar` (`id`, `remark`) VALUES ('2', 'existing')'
\end{verbatim}


\end{document}
